% Estilo v2 (2026-10-06), desenhado com a skill frontend-design e ajustado
% para ser compacto. Usado pelas disciplinas listadas em ESTILO_V2 no build.sh.
% Letras (Charter / Avenir Next / Menlo / STIX Two Math) e margens vêm do
% build.sh (-V mainfont=... -V geometry=...).
%
% Só há três caixas: definicao, exemplo e pratica. As caixas "atencao" e
% "exame" deixaram de existir (pedido do Gonçalo): o que interessa vai para o
% texto normal. Se uma aparecer, a compilação falha de propósito.
\usepackage{graphicx}
\usepackage[most]{tcolorbox}
\usepackage{xcolor}
\usepackage{fancyhdr}
\usepackage{enumitem}
\usepackage{booktabs}
\usepackage{longtable}
\usepackage{titlesec}
\usepackage{newunicodechar}
\usepackage{needspace}
\usepackage[colorlinks=true,linkcolor=Teal,urlcolor=Teal]{hyperref}

% Paleta: tinta + uma cor por função.
\definecolor{Ink}{HTML}{1F2430}
\definecolor{Teal}{HTML}{0F6B6E}     % definição e número da aula
\definecolor{Graphite}{HTML}{5B6270} % exemplo resolvido
\definecolor{Violet}{HTML}{4B3A8C}   % pratica agora
\definecolor{Paper2}{HTML}{F4F5F7}
\color{Ink}

% Menlo sem a conversão "--" -> "–" do TeX (estraga código: "->", "--").
\defaultfontfeatures{}
\setmonofont{Menlo}[Scale=0.84]
% A Charter não tem setas nem símbolos lógicos/de conjuntos: vão buscar-se
% à fonte matemática (fora do código; no código o Menlo já os tem).
\newunicodechar{→}{\ensuremath{\rightarrow}}
\newunicodechar{←}{\ensuremath{\leftarrow}}
\newunicodechar{⇒}{\ensuremath{\Rightarrow}}
\newunicodechar{↔}{\ensuremath{\leftrightarrow}}
\newunicodechar{⇔}{\ensuremath{\Leftrightarrow}}
\newunicodechar{⇏}{\ensuremath{\not\Rightarrow}}
\newunicodechar{⊆}{\ensuremath{\subseteq}}
\newunicodechar{⊂}{\ensuremath{\subset}}
\newunicodechar{∈}{\ensuremath{\in}}
\newunicodechar{∉}{\ensuremath{\notin}}
\newunicodechar{∅}{\ensuremath{\emptyset}}
\newunicodechar{∪}{\ensuremath{\cup}}
\newunicodechar{∩}{\ensuremath{\cap}}
\newunicodechar{≤}{\ensuremath{\leq}}
\newunicodechar{≥}{\ensuremath{\geq}}
\newunicodechar{≠}{\ensuremath{\neq}}
\newunicodechar{≡}{\ensuremath{\equiv}}
\newunicodechar{⊢}{\ensuremath{\vdash}}
\newunicodechar{⊨}{\ensuremath{\models}}
\newunicodechar{∧}{\ensuremath{\land}}
\newunicodechar{∨}{\ensuremath{\lor}}
\newunicodechar{¬}{\ensuremath{\lnot}}
\newunicodechar{∀}{\ensuremath{\forall}}
\newunicodechar{∃}{\ensuremath{\exists}}
\newunicodechar{×}{\ensuremath{\times}}

\linespread{1.06}
\setlength{\parskip}{4pt}
\setlength{\parindent}{0pt}
\setlist{itemsep=1pt, topsep=2pt, parsep=1pt}
% Nada de linhas soltas no fundo/topo da página, nem de hífenes a partir
% uma palavra no fim da página.
% Caixas não-quebráveis deixam espaço no fundo da página: não esticar o
% texto para o encher (era o que abria espaços enormes entre linhas).
\raggedbottom
\widowpenalty=10000
\clubpenalty=10000
\brokenpenalty=10000

% Rótulo das caixas: tipo (cor) e título, na mesma linha, sem travessão.
\newcommand{\boxlabel}[3]{{\sffamily\small\bfseries\color{#1}#2}%
  \if\relax\detokenize{#3}\relax\else\hspace{0.6em}{\sffamily\small\bfseries\color{Ink}#3}\fi\par\vspace{2pt}}
% Caixas NÃO quebráveis de propósito: uma caixa nunca se parte entre duas
% páginas. Se não couber numa página, divide-se em duas caixas.
\tcbset{base/.style={enhanced, sharp corners, frame hidden, boxrule=0pt,
  left=9pt, right=5pt, top=4pt, bottom=4pt, before skip=7pt, after skip=7pt}}
\newtcolorbox{definicao}[1][]{base, colback=white,
  borderline west={1.4pt}{0pt}{Teal}, before upper={\boxlabel{Teal}{Definição}{#1}}}
\newtcolorbox{exemplo}[1][]{base, colback=Paper2,
  before upper={\boxlabel{Graphite}{Exemplo resolvido}{#1}}}
\newtcolorbox{pratica}[1][]{base, colback=white,
  borderline west={1.4pt}{0pt}{Violet},
  borderline={0.4pt}{0pt}{Violet!35, dashed},
  before upper={\boxlabel{Violet}{Pratica agora}{#1}}}

% "Aula N": número grande à esquerda e título ao lado (div-envs.lua define
% \aulanum e passa o título sem o "Aula N ---"). Cada aula numa página nova.
% \aulanum fica vazio até à primeira aula (o título do índice usa a mesma
% caixa, sem número nem "Aula"). Nota: um título de índice com outra
% estrutura (só {\LARGE #1}) abria espaços enormes entre linhas no índice
% e noutras páginas --- causa não percebida; manter esta estrutura.
\newcommand{\aulanum}{}
\newcommand{\aulavazia}{}
\newcommand{\aulahead}[1]{%
  \begin{minipage}[b]{2.4em}{\fontsize{44}{40}\selectfont\bfseries\color{Teal}\aulanum}\end{minipage}%
  \hspace{0.6em}\begin{minipage}[b]{\dimexpr\linewidth-3em}\raggedright
  \ifx\aulanum\aulavazia\else{\small\bfseries\color{Teal}Aula \aulanum}\\[1pt]\fi
  {\LARGE\bfseries #1}\end{minipage}}
% O pandoc não numera secções, por isso o titlesec usa o formato "numberless".
\titleformat{\section}[block]{\sffamily\color{Ink}}{}{0pt}{\aulahead}
\titleformat{name=\section,numberless}[block]{\sffamily\color{Ink}}{}{0pt}{\aulahead}
\titlespacing*{\section}{0pt}{0pt}{14pt}
\titleformat{\subsection}{\sffamily\large\bfseries\color{Ink}\raggedright}{}{0pt}{}
\titlespacing*{\subsection}{0pt}{12pt}{3pt}
\titleformat{\subsubsection}{\sffamily\normalsize\bfseries\color{Graphite}\raggedright}{}{0pt}{}
\titlespacing*{\subsubsection}{0pt}{8pt}{2pt}
% A página nova antes de cada aula é posta pelo div-envs.lua (\clearpage
% antes do \section), para o título e o índice ficarem na primeira página.

\pagestyle{fancy}
\fancyhf{}
\fancyhead[L]{\footnotesize\sffamily\color{Graphite}\nouppercase{\leftmark}}
\fancyhead[R]{\footnotesize\sffamily\color{Graphite}\thepage}
\renewcommand{\headrulewidth}{0pt}
\renewcommand{\sectionmark}[1]{\markboth{#1}{}}

% --- Deduções em notação de Fitch (geradas por fitch.lua a partir de ```{.fitch}) ---
\newcommand{\fitchbar}{\rule[-0.4em]{0.6pt}{1.45em}\hspace{0.55em}}
\newcommand{\fitchbarstart}{\rule[-0.4em]{0.6pt}{1.2em}\hspace{0.55em}}
\newcommand{\fitchhyp}[1]{\hspace{-0.5em}\underline{\hspace{0.5em}#1\hspace{0.9em}}}
\newenvironment{fitchproof}{%
  \par\smallskip\begingroup\centering\renewcommand{\arraystretch}{1}%
  \begin{tabular}{@{}r@{\hspace{1em}}l@{\hspace{2.2em}}l@{}}}%
  {\end{tabular}\par\endgroup\smallskip}

% Título do documento compacto, na mesma página do índice.
\makeatletter
\renewcommand{\maketitle}{{\sffamily\raggedright
  {\LARGE\bfseries\color{Ink}\@title\par}\vspace{2pt}
  {\small\color{Graphite}\@author\quad\@date\par}}\vspace{10pt}\thispagestyle{empty}}
\makeatother
% Índice compacto (cabe numa página). Com o tocloft as listas dentro das
% caixas ganhavam espaços enormes entre itens: usar o titletoc.
\usepackage{titletoc}
\titlecontents{section}[0pt]{\addvspace{4pt}\sffamily\bfseries\small}{}{}{\hfill\contentspage}
\titlecontents{subsection}[1.5em]{\small}{}{}{\titlerule*[0.6em]{.}\contentspage}
